\subsection{A necessary vertical fibre over the missing origin}

\begin{proposition}[B.4.26; Proved; every origin component is a threefold]
\label{push:B426}
Suppose the actual partial family of B.4.19 has an
proper analytic completion $f:\mathcal Z\to\Delta_q\times\Delta_r$
with smooth pure four-dimensional total space, with
$r\circ f$ flat, and with its prescribed identification
over the punctured bidisc. The map to this fixed bidisc
is assumed holomorphic. Then every irreducible component
of $f^{-1}(0,0)$ has dimension three. In particular,
there is no equidimensional completion, and no completion
having even one two-dimensional component of the origin
fibre.

The conclusion still permits the vertical threefold
components allowed in B.4. It is not an assertion that
an arbitrary $r$-proper filling initially has a holomorphic
extension of $q$; if $q$ is only meromorphic, a graph
modification must first be supplied and its resulting
geometry retained.
\end{proposition}

\begin{proof}
We first prove the finite-multisection consequence of
the actual log-torsor obstruction B.4.24, without
assuming that the unknown cap is itself a log torsor.
Make the normalized finite base change $q=Q^{12}$.
The resulting space $\mathcal Z'$ is normal and
pure four-dimensional. Its restrictions to the two
punctured axes have the explicit normalized fan
presentations used in B.4.24. Write
\[
 B=\Delta_Q\times\Delta_r,\quad
 B_q=\Delta_Q^*\times\Delta_r,\quad
 B_r=\Delta_Q\times\Delta_r^*,\quad
 U=B_q\cap B_r.
\]

Assume there is a finite surjective map
$p:V\to B$ of positive generic degree $d$, together
with a holomorphic map $V\to\mathcal Z'$ over $B$.
We may normalize $V$ and retain any collection of
components dominating $B$; its degree remains a
positive integer. A normal surface is Cohen--Macaulay;
finite over the smooth bidisc, it is flat. Thus the
finite holomorphic trace and the degree-$d$ sum with
multiplicities are available also at branch points.
On the ordinary locus, sum the
$d$ points of this multisection, with multiplicity,
in each of the two actual axial group origins of
B.4.24. This gives two local expressions for the
trace in the Jacobian group. We show that they
extend as log-torus sections on $B_q$ and $B_r$.

\emph{Ramification away from the axes.}
Over the ordinary locus the torus family is a smooth
holomorphic group family. Near a branch point, choose
small neighbourhoods of the finitely many preimages
in $V$ and local Lie lifts of their maps to the torus.
The trace of these holomorphic functions under the
finite map is holomorphic, and its exponential is
the ordinary sum. Changes of lifts add periods,
which disappear in the group quotient. Thus the
ordinary trace extends across the ordinary branch
locus. Equivalently this is the usual addition map
on the relative symmetric power, applied to the
finite multisection.

\emph{The explicit norm at an axis.}
After shrinking $B$, every component of the finite
cover's branch divisor which is not an axis meets
each punctured axis nowhere: its intersection with
either axis is an isolated analytic germ at the
origin. At a point of a punctured axis choose
coordinates $(x,y)$, with $x$ the normal coordinate
to that axis. Locally the normal finite cover is
a disjoint union of covers
\[
 x=t^e,\qquad y=y,
 \tag{B.4.26.1}
\]
including $e=1$ sheets. Indeed the complement of
the axis in a small polydisc has cyclic fundamental
group; each connected unramified finite cover is
the displayed root cover there, and its normal
finite extension is unique. No assumption that the
whole finite cover is a coordinate-root cover is
being made.

The map from such a root polydisc into the
prescribed fan model gives a section of its
associated log-torus sheaf. It has a local bounded
multiplicative log lift. This can also be seen
directly in a toric chart: the pulled-back characters
are meromorphic functions with zeros and poles
supported on $t=0$, hence are $t^k u(t,y)$ with
$k\in\Z$ and $u$ a holomorphic unit. The quotient
by the period lattice causes only the period
ambiguity already retained in B.4.24.

The multiplicative norm of a character on (B.4.26.1)
is explicitly
\[
 \prod_{j=0}^{e-1}(\zeta_e^j t)^k
              u(\zeta_e^j t,y)
 =(-1)^{k(e-1)}x^k
              \prod_{j=0}^{e-1}u(\zeta_e^j t,y).
 \tag{B.4.26.2}
\]
The last product is a holomorphic unit in $(x,y)$.
Thus the norm has an integral ghost vector and
is an actual log character on the base. It is
bounded in the required sense: on the $r$ axis
the first character has zero normal order and
the second is unrestricted; on the $Q$ axis
both are unrestricted. These are precisely the
bounded-character conditions for (B.4.24.1).
Multiplication of lifts corresponds to addition
in the log torus. Changing a lift by a period
changes its norm by an integral period, so the
resulting quotient section is independent of
the lift. Summing over the finitely many covers
(B.4.26.1), whose degrees add to $d$, proves the
claimed extension along that axis.

The local extensions agree on overlaps because
they agree on the ordinary locus. To see uniqueness
without a compactification assumption, lift locally
to log characters; agreement there is an equality
up to a locally constant $Y$-period, and that
meromorphic equality extends across the axis.
We have therefore obtained actual sections
$h_q$ of $\mathcal A$ on $B_q$ and $h_r$ on
$B_r$. The change between the two origins for
each point is $F$, so on $U$
\[
 h_q-h_r=dF
 \tag{B.4.26.3}
\]
up to the harmless overall choice of subtraction
convention. The ghost proof of B.4.24 forbids
(B.4.26.3): its $r$-loop $Y$-component is
$d\gamma\ne0$. Thus no such finite multisection
can exist.

\emph{A two-dimensional fibre germ would supply one.}
Suppose an origin point of $\mathcal Z$ has local
fibre dimension two. Finite base change and
normalization are finite maps and preserve that
local fibre dimension at points lying above it.
At a resulting point of $\mathcal Z'$, choose
two generic local linear sections in an analytic
embedding. Their intersection with the fibre is
isolated. Intersecting with $\mathcal Z'$ gives
a two-dimensional analytic germ whose map to
$B$ is finite. This is the local finite-projection
criterion from Weierstrass preparation; see also
\citet[Chapter II, Sections 7--8]{Demailly2012E}.
Choose a component through the point. Its finite
image has dimension two and is therefore the
whole base germ. A finite representative of this
germ, after shrinking the bidisc, is finite and
surjective over the entire smaller bidisc. Its
normalization is still finite and maps to
$\mathcal Z'$, giving exactly the forbidden
multisection of positive degree. Properness of
this finite representative prevents loss of
sheets at its boundary.

Every fibre component of a dominant map from a
pure fourfold to a smooth surface has dimension
at least two. If a two-dimensional irreducible
component existed, a generic point on it would
have local fibre dimension two, which we have
excluded. Flatness over $r$ makes $r=0$ a
Cartier divisor of pure dimension three, so
the origin fibre has no component of dimension
greater than three. All its irreducible
components consequently have dimension three,
as asserted.
\end{proof}
